\section{Refleksja filozoficzna w świetle hipotezy Sapira-Whorfa}
\label{sec:philosophy}

\subsection{Hipoteza i stan sporu}

Humboldt pisał, że budowa języka wpływa na rozwój umysłowy ludzkości \cite{humboldt1836}.
Sapir twierdził w roku 1929, że człowiek jest zdany na język społeczności, a obraz
świata opiera się na nawykach językowych grupy \cite{sapir1929}. Whorf sformułował w roku
1940 zasadę relatywizmu językowego: te same dane fizyczne prowadzą obserwatorów do tego
samego obrazu świata tylko przy podobnym zapleczu językowym \cite{whorf1940,whorf1956}.
Nazwa \pkq{hipoteza Sapira-Whorfa} pochodzi od następców obu autorów.

Literatura rozróżnia 2 wersje hipotezy. Wersja mocna głosi, że język wyznacza myślenie.
Wersja słaba głosi, że język wpływa na myślenie. Kay i Kempton wykryli wpływ nazw barw na
ocenę podobieństwa barw \cite{kay1984}. Winawer i współautorzy pokazali, że osoby
mówiące po rosyjsku szybciej rozróżniają 2 odcienie niebieskiego, gdy odcienie mają
w języku rosyjskim różne nazwy \cite{winawer2007}. Tom pod redakcją Gumperza i Levinsona
porządkuje spór \cite{gumperz1996}. Wersja mocna nie ma potwierdzenia. Wersja słaba ma
potwierdzenie w wybranych dziedzinach.

\subsection{Tradycja polska}

Polska filozofia ma własny wariant tezy. Ajdukiewicz ogłosił w roku 1934 radykalny
konwencjonalizm: obraz świata zależy od aparatury pojęciowej, czyli od wyboru języka
\cite{ajdukiewicz1934}. Schaff poświęcił hipotezie Sapira-Whorfa monografię
\pkq{Język a poznanie} \cite{schaff1964}. Bartmiński rozwinął pojęcie językowego obrazu
świata \cite{bartminski2006}. Wierzbicka szuka pojęć prostych, które są wspólne
językom \cite{wierzbicka1996}.

Szkoła lwowsko-warszawska wiązała jasność języka z jasnością myślenia. Twardowski
twierdził, że niejasny styl ma źródło w niejasnym myśleniu autora \cite{twardowski1919}.
Tarski wykazał, że ścisła definicja prawdy wymaga języka o określonej strukturze
\cite{tarski1933}. PK korzysta z obu myśli. Reguły PK działają w kierunku przeciwnym do
diagnozy Twardowskiego: wymóg jasnego zapisu wymusza jasną decyzję. Wynik Tarskiego
wyznacza granicę: PK 0.1 nie ma struktury formalnej, a język \PKz{} ma.

\subsection{Ideał jednego znaczenia}

PK należy do tradycji języka idealnego. Leibniz szukał alfabetu myśli, w którym spór
zamienia się w rachunek \cite{leibniz1666}. Frege zbudował pismo pojęciowe i porównał
stosunek tego pisma do języka potocznego ze stosunkiem mikroskopu do oka
\cite{frege1879}. Wittgenstein pisze w \pkq{Traktacie logiczno-filozoficznym}, że język
potoczny przesłania formę myśli (teza 4.002) i że każdą myśl da się wypowiedzieć jasno
(teza 4.116) \cite{wittgenstein1921}. Teza 5.6 wiąże granice języka z granicami świata.
Carnap dodał zasadę tolerancji: zadaniem logika jest dochodzenie do konwencji, a nie
ustanawianie zakazów \cite{carnap1934}. Quine opisał regimentację, czyli parafrazę zdań
potocznych na notację kanoniczną \cite{quine1960}. Montague odrzucił istotną różnicę
teoretyczną między językiem naturalnym a językiem formalnym \cite{montague1970}. Brown opisał
w roku 1960 język logiczny Loglan \cite{brown1960}; Loglan jest nowym językiem, a PK jest
podzbiorem polszczyzny.

PK realizuje ten ideał skromnie. PK nie zastępuje polszczyzny notacją. PK wybiera
z polszczyzny zdania, które czytelnik odczytuje w 1 sposób, i dokłada konwencje. Zasada
tolerancji Carnapa opisuje status PK trafnie: PK jest konwencją, którą autor przyjmuje
dla 1 rodzaju tekstu.

Tradycja ma także krytyków. Późny Wittgenstein twierdzi, że reguła nie wyznacza własnego
zastosowania; zastosowanie ustala praktyka wspólnoty \cite{wittgenstein1953}. Reguła E6
standardu PK mówi to samo językiem inżyniera: konwencja działa tylko dla czytelnika,
który zna konwencję. Grice umieścił wśród maksym sposobu 2 zalecenia: unikaj
wieloznaczności oraz mów zwięźle \cite{grice1975}. PK rozstrzyga konflikt tych maksym na
korzyść jednoznaczności; pomiar P3 podaje cenę. Piantadosi, Tily i Gibson pokazali, że
niejednoznaczność oszczędza wysiłek, gdy kontekst rozstrzyga znaczenie
\cite{piantadosi2012}. PK jest więc uzasadniona tam, gdzie kontekstu brakuje albo gdzie
błąd odczytania kosztuje więcej niż dłuższy tekst.

\subsection{PK jako zestaw kategorii obowiązkowych}

Jakobson, za Boasem, ujął różnicę między językami jednym zdaniem: \pkq{Languages differ
essentially in what they must convey and not in what they may convey}
\cite{jakobson1959}. Polszczyzna każe podać rodzaj rzeczownika i aspekt czasownika.
Polszczyzna pozwala pominąć wykonawcę, ilość, granicę przedziału i źródło wiedzy.

PK zmienia listę kategorii obowiązkowych dla 1 rejestru. Autor tekstu w PK podaje 5 danych.
\begin{itemize}[leftmargin=1.6em,itemsep=0pt]
  \item Wykonawca działania (reguła B1).
  \item Ilość albo nazwa obiektu (reguła A2).
  \item Strona granicy przedziału (reguła D2).
  \item Chwila w kalendarzu (reguła D3).
  \item Źródło wiedzy (moduł F).
\end{itemize}
Moduł F jest odpowiednikiem ewidencjalności. Aikhenvald opisała języki, w których każde zdanie wskazuje źródło
informacji: mówiący widział zdarzenie, słyszał o zdarzeniu albo wnioskuje
\cite{aikhenvald2004}. PK wprowadza tę kategorię do polszczyzny konwencją, a nie gramatyką.

Slobin nazwał wpływ kategorii obowiązkowych \pkq{myśleniem dla mówienia}: mówiący kieruje
uwagę na cechy, które gramatyka każe wyrazić \cite{slobin1996}. Praca stawia analogiczną
hipotezę dla pisania. Autor, który pisze w PK, rozstrzyga fakty, które zwykła polszczyzna
pozwala zostawić bez rozstrzygnięcia. Gdy autor faktu nie zna, standard każe wstawić
znacznik \pkq{[DO USTALENIA]}. Przykład ze standardu pokazuje skalę zjawiska: 3 zdania
instrukcji dają 3 znaczniki braku. PK ujawnia więc niewiedzę autora przed wykonaniem
instrukcji. Hipoteza jest słabą wersją tezy Sapira-Whorfa dla 1 rejestru. Praca tej
hipotezy nie testuje; test wymaga badania z udziałem ludzi.

\subsection{Hipoteza Sapira-Whorfa a model językowy}
\label{sec:swllm}

Model językowy jest funkcją: ciąg tokenów wyznacza rozkład prawdopodobieństwa kolejnego
tokenu. Praca rozdziela pytanie o hipotezę Sapira-Whorfa dla modelu na 3 tezy o różnym
statusie.

\paragraph{Teza A: zależność od tekstu.}
Zadanie dociera do modelu wyłącznie jako tekst. Odpowiedź modelu zależy więc od
sformułowania zadania. Teza A zachodzi z samej budowy modelu i nie wymaga pomiaru.
Dla człowieka mocna wersja hipotezy jest sporna, bo człowiek poznaje świat także poza
językiem. Model językowy takiej drogi nie ma: granice tekstu wejściowego są granicami
zadania. Bender i Koller ostrzegają przed wnioskiem, że system wyuczony na samej formie
języka ma dostęp do znaczenia \cite{bender2020}. Praca stosuje tę ostrożność: słowo
\pkq{odczytanie} oznacza w pomiarze P4 rozkład odpowiedzi modelu, a nie stan umysłu.

\paragraph{Teza B: efekt rejestru.}
Ta sama treść zapisana w PK i w zwykłej polszczyźnie daje różne rozkłady odpowiedzi.
Teza B jest empiryczna; wynik podaje sekcja \ref{sec:probe}.

\paragraph{Teza C: efekt nawyku.}
Model uczy się na tekstach w zwykłej polszczyźnie, więc nie zna konwencji PK. Konwencja
działa na model dopiero wtedy, gdy tekst zawiera legendę. Teza C przenosi regułę E6 na
model językowy; wynik podaje sekcja \ref{sec:probe}.

Literatura dostarcza 2 obserwacji o języku wejścia. Wendler i współautorzy pokazali, że
w modelach rodziny Llama-2 reprezentacje warstw pośrednich leżą bliżej języka angielskiego
niż języka zapytania \cite{wendler2024}. Kim i współautorzy zmierzyli 26 języków w zadaniach
wyszukiwania w tekście o długości 64 tys. i 128 tys. tokenów \cite{kim2025}. Język polski
zajął miejsce 1, a język angielski miejsce 6. Wynik dotyczy 7 zadań tego testu
i nie dowodzi ogólnej przewagi polszczyzny. Oba wyniki mówią o języku naturalnym jako
całości. Niniejsza praca bada węższe pytanie: rejestr wewnątrz 1 języka.

Hipoteza ma więc dla modelu inny kształt niż dla człowieka. U człowieka spór dotyczy
trwałego wpływu języka ojczystego na poznanie. U modelu wpływ tekstu wejściowego jest
doraźny i całkowity, a ciekawe pytanie brzmi: które cechy tekstu zmieniają rozkład
odpowiedzi.

\subsection{Granice}
\label{sec:limits}

\paragraph{Nowomowa.}
Orwell opisał język zaprojektowany po to, żeby zawęzić zakres myśli
\cite{orwell1949}. Model z sekcji \ref{sec:model} pokazuje różnicę między nowomową a PK.
Nowomowa zmniejsza zbiór znaczeń, które użytkownik potrafi wyrazić, czyli obraz relacji
interpretacji w zbiorze $M$. PK zmniejsza zbiory odczytań pojedynczych zdań. Dla każdego
ostrego odczytania zdania zwykłego standard PK wskazuje osobne zdanie PK; tabela
w dodatku \ref{app:items} zawiera \pkItemsPkSentences{} takich zdań. Nowomowa odbiera
rozróżnienia, a PK wymusza rozróżnienia. PK jest także dobrowolna i obejmuje 1 rejestr.

\paragraph{Nieostrość celowa.}
PK usuwa znaczenia nieostre, a nieostrość bywa potrzebna. Ustawodawca pisze
\pkq{niezwłocznie} albo \pkq{należyta staranność}, żeby zostawić sądowi ocenę przypadku.
Pomiar P5 pokazuje, że akty prawne używają takich słów stale. Standard PK zastrzega, że
umowa w PK nadal wymaga oceny prawnika. PK nie nadaje się także do literatury, reklamy
i rozmowy; standard wyłącza te rodzaje tekstu.

\paragraph{Konwencja wymaga wspólnoty.}
Jednoznaczność zdania PK nie jest własnością samego ciągu słów. Zdanie \pkq{Ciśnienie
zmienia przepływ.} jest tym samym ciągiem słów w PK i poza PK. Jednoznaczność powstaje
dopiero dla czytelnika, który zna regułę B2. PK jest więc umową między autorem
a czytelnikiem. Rysunek \ref{fig:timeline} umieszcza PK na tle 3 tradycji, z których
standard korzysta.

\rysunek{timeline}{Historia idei: wybrane prace 3 tradycji, które łączy PK.}{fig:timeline}

\section{Zagrożenia dla trafności}
\label{sec:threats}

\paragraph{Trafność konstrukcji.}
Sonda wymusza odpowiedź \pkq{Tak} albo \pkq{Nie}. Model nie ma opcji \pkq{zdanie jest
niejednoznaczne}. Pytanie sondy wskazuje strefę szarą, więc sonda mierzy wybór odczytania,
a nie wykrycie niejednoznaczności. Zdania pomiaru stoją bez kontekstu. W tekście kontekst
rozstrzyga część niejednoznaczności, więc pomiar zawyża skalę problemu zdań zwykłych.

\paragraph{Trafność wewnętrzna.}
Ten sam autor napisał standard, pozycje testowe, klucz odpowiedzi i sprawdzarkę.
Autor napisał \pkItemsAuthor{} z \pkItems{} pozycji po zaprojektowaniu pomiaru.
Klasyfikację reguł i ocenę precyzji wykonał 1 oceniający. Praca nie ma oceny językoznawcy.
Autor ustalił próg kontroli kotwic, czyli 80 \%, po obejrzeniu wyniku pierwszego modelu.
Pozycja K2 ma sporny klucz (sekcja \ref{sec:probe}).

\paragraph{Trafność zewnętrzna.}
Modele lokalne mają od \pkLocalParamMin{} do \pkLocalParamMax{} mld parametrów włącznie; tor Claude
obejmuje 2 modele.
Wynik nie przenosi się wprost na inne modele. Pozycje opisują instalacje, terminy
i regulaminy; inne dziedziny wymagają osobnego pomiaru. Praca nie zawiera badania
z udziałem ludzi.

\paragraph{Wnioskowanie statystyczne.}
Pomiar P4 ma \pkItemsRewrite{} pozycji typu przeredagowanie, \pkItemsConv{} pozycje typu
konwencja i \pkItemsGap{} pozycje typu luka. Wynik dla konwencji i dla luk jest opisem
przypadków, a nie estymacją. Test znaków nie ma poprawki na liczbę modeli. Tor Claude
ma 2 próby na sondę.

\paragraph{Model formalny.}
Twierdzenie \ref{thm:product} zakłada niezależne miejsca wyboru. Twierdzenie
\ref{thm:questions} wymaga rozkładu zamierzonych odczytań, którego praca nie zna.
Twierdzenia \ref{thm:horn} i \ref{thm:np} dotyczą języka \PKz, a nie pełnej PK 0.1.
Gramatyki pomiaru P2 modelują 1 konstrukcję. Twierdzenie \ref{thm:breakeven} zakłada
niezależne próby; model językowy z temperaturą 0 powtarza ten sam błąd, więc wzór
opisuje zbiór zadań, a nie powtórzenia 1 zadania.

\paragraph{Sprawdzarka.}
Sprawdzarka jest przybliżona. Precyzja pochodzi z próby \pkPrecN{} trafień; czułość jest
nieznana. Sprawdzarka nie obejmuje reguł, które wymagają składni albo znaczenia.
